\documentclass[../../../include/open-logic-section]{subfiles}

\begin{document}

\olfileid{lam}{int}{pr}
\olsection{$\lambda$-可定义函数在原始递归下封闭}

% Part: lambda-calculus
% Chapter: introduction
% Section: primitive-recursion

至于原始递归，我们终于需要花点功夫了。我们必须分阶段进行。如前所述，假设我们已经有分别$\lambd$-定义函数$g$和~$h$的项$G$和~$H$，我们需要找到一个项$F$来$\lambd$-定义由
\begin{align*}
f(0, \vec z) & = g(\vec z) \\
f(x+1, \vec z) & = h(z, f(x,\vec z), \vec z).
\end{align*}
定义的函数~$f$。
因此，一般地，给定$\lambda$项$G'$和~$H'$，只需找到一个项~$F$使得
\begin{align*}
F(\num{0}, \vec z) & \equiv G(\vec z) \\
F(\overline{n+1}, \vec z) & \equiv H(\num{n}, F(\num{n}, \vec z), \vec z)
\end{align*}
对每个自然数~$n$成立即可；既然$G'$和~$H'$分别$\lambd$-定义了$g$和~$h$，这就意味着只要将数字$\num{\vec m}$代入$\vec z$，$F(\num{n+1}, \num{\vec m})$就会归约到正确的答案。

为此，只需找到一个满足
\begin{align*}
F(\num 0) & \equiv G \\
F(\overline {n+1}) & \equiv H(\num{n},F(\num{n}))
\intertext{的项~$F$即可，其中对每个自然数$n$，有}
G & = \lambd[\vec z][G'(\vec z)] \text{ 且}\\
H(u,v) & = \lambd[\vec z][H'(u,v(u,\vec z),\vec z)].
\end{align*}
换言之，借助$\lambda$技巧，我们无需顾虑额外的参数$\vec z$——它们直接被吸收进了$\lambda$记法中。

在定义项$F$之前，我们需要一个处理有序对的机制。这由下一个引理提供。

\begin{lem}
存在一个λ-项 $D$，使得对每一对λ-项 $M$ 与~$N$，$D(M,N)(\num{0}) \red M$ 且 $D(M,N)(\num{1}) \red N$。
\end{lem}

\begin{proof}
首先，定义λ-项 $K$ 如下：
\[
K(y) = \lambd[x][y].
\]
换句话说，$K$ 是项 $\lambd[y][\lambd[x][y]]$。换个角度看，对每个 $M$，$K(M)$ 都是一个常函数，无论输入什么都会返回~$M$。

现在，通过 $D(x,y,z) = z (K(y))x$ 定义 $D(x,y,z)$。于是我们有
\begin{align*}
D(M,N,\num 0) & \red \num 0 (K(N)) M \red M \text{ 且 }\\
D(M,N,\num 1) & \red \num 1 (K(N)) M \red K(N) M \red N,
\end{align*}
如所求。
\end{proof}

其核心想法是，$D(M,N)$ 表示有序对 $\tuple{M, N}$，并且若假定 $P$ 表示这样一个有序对，则 $P(\num 0)$ 和 $P(\num 1)$ 便分别表示左投影和右投影，即 $(P)_0$ 和 $(P)_1$。我们将采用后一种记法。

\begin{lem}
λ-可定义函数在原始递归下封闭。
\end{lem}

\begin{proof}
我们需要证明，对于任意给定的项 $G$ 和 $H$，都能找到一个项 $F$，使得对每个自然数~$n$，
\begin{align*}
F(\num{0}) & \equiv G \\
F(\num{n+1}) & \equiv H(\num{n}, F(\num{n}))
\end{align*}。
大致思路是，以数字作为迭代器来计算\emph{有序对}的序列
\[
\tuple{\num 0, F(\num 0)}, \tuple{\num 1, F(\num 1)}, \dots,
\]
注意，第一个有序对就是 $\tuple{\num 0, G}$。给定一个有序对 $\tuple{\num{n}, F(\num{n})}$，下一个有序对 $\tuple{\num{n+1}, F(\num{n+1})}$ 应当与 $\tuple{\num{n+1}, H(\num{n}, F(\num{n}))}$ 等价。我们将设计一个实现这一单步转移的λ-项~$T$。

细节如下。定义 $T(u)$ 为：
\[
T(u) = \tuple{S((u)_0), H((u)_0,(u)_1)}.
\]
现在容易验证，对任意自然数~$n$，
\[
T(\tuple{\num{n}, M}) \red \tuple{\num{n+1}, H(\num{n}, M)}.
\]
如上所述，给定 $G$ 和 $H$，通过下式定义~$F(u)$：
\[
F(u) = (u(T,\tuple{\num 0, G}))_1.
\]
换言之，以~$\num{n}$ 为输入，$F$ 将 $T$ 在 $\tuple{\num 0, G}$ 上迭代 $n$ 次，然后返回第二个分量。首先，我们有
\begin{enumerate}
\item $\num{0} (T, \tuple{\num 0, G}) \equiv \tuple{\num{0}, G}$
\item $F(\num{0}) \equiv G$
\end{enumerate}
通过对~$n$ 归纳，可证对每个自然数均有以下结论：
\begin{enumerate}
\item $\num{n+1}(T, \tuple{\num{0}, G}) \equiv \tuple{\num{n+1}, F(\num{n+1})}$
\item $F(\num{n+1}) \equiv H(\num{n}, F(\num{n}))$
\end{enumerate}
对于第二条，我们有
\begin{align*}
F(\num{n+1}) & \red (\num{n+1}(T, \tuple{\num 0, G}))_1 \\
& \equiv (T(\num{n} (T, \tuple{\num 0, G})))_1 \\
& \equiv (T(\tuple{\num{n}, F(\num{n})}))_1 \\
& \equiv (\tuple{\num{n+1}, H(\num{n}, F(\num{n}))})_1 \\
& \equiv H(\num{n}, F(\num{n})).
\end{align*}
此处我们在倒数第二行用到了归纳假设。对于第一条，我们有
\begin{align*}
\num{n+1} (T, \tuple{\num 0, G}) &
\equiv T(\num{n} (T, \tuple{\num 0, G})) \\
& \equiv T( \tuple{\num{n}, F(\num{n})}) \\
& \equiv \tuple{\num{n+1}, H(\num{n}, F(\num{n}))} \\
& \equiv \tuple{\num{n+1}, F(\num{n+1})}.
\end{align*}
此处我们在最后一行用到了第二条。因此，我们证明了 $F(\num 0) \equiv G$ 且对每个 $n$ 有 $F(\num {n+1}) \equiv H(\num
n, F(\num{n}))$，而这正是我们所需要的。
\end{proof}

\end{document}